\documentclass[12pt]{article}

\usepackage[a4paper, margin=1in]{geometry}
\usepackage{enumitem}
\usepackage{multicol}
\usepackage{amsmath}
\usepackage{times}
\usepackage{fancyhdr}
\usepackage{tabularx}
\usepackage{needspace}
\usepackage{xcolor}

\pagestyle{fancy}
\fancyhf{}
\renewcommand{\headrulewidth}{0pt}
\renewcommand{\footrulewidth}{0pt}
\fancyfoot[R]{Page \thepage}
\fancyfoot[L]{Twiga Generated Practice Exam}

\newcommand{\leftheading}[1]{
    \noindent{\fontsize{12}{14}\selectfont\bfseries #1}\par
}

\newcommand{\centerheading}[1]{
    \vspace{0.6em}
    \begin{center}
        {\fontsize{13}{15}\selectfont\bfseries #1}
    \end{center}
    \vspace{0.15em}
}

\newenvironment{questionblock}[1][10]{
    \par\Needspace{#1\baselineskip}
}{
    \par
}

% defining the solution text color and layout commands
\definecolor{solutionred}{RGB}{170,0,0}
\newcommand{\solutionheading}[1]{
    \par{\color{solutionred}\textbf{#1}}\par
}
\newcommand{\solutionline}[1]{
    \par{\color{solutionred}#1}\par
}
\newenvironment{solutionblock}{
    \par\begingroup\color{solutionred}
}{
    \par\endgroup
}

\begin{document}

\begin{center}
  {\fontsize{14}{16}\selectfont\bfseries THE UNITED REPUBLIC OF TANZANIA\par}
  \vspace{0.2em}
  {\fontsize{12}{14}\selectfont PRESIDENT’S OFFICE\par}
  {\fontsize{12}{14}\selectfont REGIONAL ADMINISTRATION AND LOCAL GOVERNMENT\par}
  \vspace{0.35em}
  {\fontsize{14}{16}\selectfont\bfseries GENERATED PRACTICE EXAM\par}
  {\fontsize{14}{16}\selectfont\bfseries CHEMISTRY\par}
\end{center}
\vspace{0.6em}
\noindent{\fontsize{12}{14}\selectfont Time: 3:00 Hrs\hfill YEAR: 2026\par}
\vspace{0.8em}

\vspace{-1.0em}
\noindent\rule{\textwidth}{0.4pt}
\vspace{-1.0em}

\leftheading{INSTRUCTIONS}
\begin{enumerate}[itemsep=0pt, topsep=2pt, parsep=0pt, partopsep=0pt]
  \item This paper consists of three sections A, B and C with a total of 18 questions.
  \item Answer all questions in section A and B and one (1) question from section C.
  \item Calculators and cellular phones are not allowed in the examination room.
  \item Write your Examination number on every page of your answer sheets/booklet(s).
\end{enumerate}

% section A
\centerheading{SECTION A (15 Marks)}
\vspace{-2.0em}
\begin{center}
  Answer all questions in this section.
\end{center}

\begin{enumerate}
  \item For each of the following items, choose the correct answer among the given alternatives and write its letter beside the item number provided. \textbf{(10 Marks)}

  \begin{enumerate}[label=(\roman*)]

    \begin{questionblock}
      \item Which of the following particles is found in the nucleus of an atom?
      \begin{enumerate}[label=\Alph*.]
        \item Electron
        \item Proton
        \item Neutron and Electron
        \item Proton and Electron
        \item Proton and Neutron
      \end{enumerate}
      \vspace{0.8em}
    \end{questionblock}

    \begin{questionblock}
      \item Which of the following elements is a noble gas?
      \begin{enumerate}[label=\Alph*.]
        \item Neon
        \item Copper
        \item Carbon
        \item Oxygen
        \item Nitrogen
      \end{enumerate}
      \vspace{0.8em}
    \end{questionblock}

    \begin{questionblock}
      \item What type of chemical bond is typically formed between a metal and a non-metal?
      \begin{enumerate}[label=\Alph*.]
        \item Covalent bond
        \item Ionic bond
        \item Hydrogen bond
        \item Metallic bond
        \item Van der Waals bond
      \end{enumerate}
      \vspace{0.8em}
    \end{questionblock}

    \begin{questionblock}
      \item What is the main characteristic that distinguishes the atomic number of an element from its mass number?
      \begin{enumerate}[label=\Alph*.]
        \item The atomic number is the total number of protons and neutrons, while the mass number is the number of protons.
        \item The atomic number is the number of protons, while the mass number is the total number of protons and neutrons.
        \item The atomic number is the number of neutrons, while the mass number is the total number of protons and electrons.
        \item The atomic number is the total number of electrons and neutrons, while the mass number is the number of protons.
        \item The atomic number is the number of electrons, while the mass number is the total number of protons, neutrons, and electrons.
      \end{enumerate}
      \vspace{0.8em}
    \end{questionblock}

    \begin{questionblock}
      \item Which of the following elements is likely to be found in group 11 of the periodic table?
      \begin{enumerate}[label=\Alph*.]
        \item Sodium
        \item Copper
        \item Neon
        \item Oxygen
        \item Nitrogen
      \end{enumerate}
      \vspace{0.8em}
    \end{questionblock}

    \begin{questionblock}
      \item What type of bond is formed when two atoms share one or more pairs of electrons?
      \begin{enumerate}[label=\Alph*.]
        \item Ionic bond
        \item Covalent bond
        \item Hydrogen bond
        \item Electrostatic bond
        \item Metallic bond
      \end{enumerate}
      \vspace{0.8em}
    \end{questionblock}

    \begin{questionblock}
      \item The atomic number of an element is defined as the number of
      \begin{enumerate}[label=\Alph*.]
        \item neutrons in the nucleus
        \item protons in the nucleus
        \item electrons in the outermost shell
        \item neutrons and protons in the nucleus
        \item protons and electrons in an atom
      \end{enumerate}
      \vspace{0.8em}
    \end{questionblock}

    \begin{questionblock}
      \item Which of the following elements is a transition metal?
      \begin{enumerate}[label=\Alph*.]
        \item Copper
        \item Carbon
        \item Nitrogen
        \item Oxygen
        \item Neon
      \end{enumerate}
      \vspace{0.8em}
    \end{questionblock}

    \begin{questionblock}
      \item What is the primary force responsible for holding atoms together in a chemical bond?
      \begin{enumerate}[label=\Alph*.]
        \item Gravitational force
        \item Electromagnetic force
        \item Nuclear force
        \item Frictional force
        \item Intermolecular force
      \end{enumerate}
      \vspace{0.8em}
    \end{questionblock}

    \begin{questionblock}
      \item The maximum number of electrons that can occupy the third energy level (n = 3) of an atom is
      \begin{enumerate}[label=\Alph*.]
        \item 8
        \item 18
        \item 2
        \item 32
        \item 50
      \end{enumerate}
      \vspace{0.8em}
    \end{questionblock}

  \end{enumerate}

\end{enumerate}

% section B
\Needspace{21\baselineskip}
\centerheading{SECTION B (70 Marks)}
\vspace{-2.0em}
\begin{center}
  Answer all questions in this section.
\end{center}
\vspace{-2.0em}

\begin{enumerate}[start=3]
  \begin{questionblock}[15]
    \item \textbf{(14 Marks)}

    (a) Describe the atomic structure components and their arrangement. (6)
    \begin{enumerate}[label=(\roman*)]
      \item Identify the nucleus contents. (2)
      \item Describe electron arrangement around the nucleus. (2)
      \item State the overall atomic charge. (2)
    \end{enumerate}

    (b) Explain the relation between atomic number, mass number, and atomic structure. (8)
    \begin{enumerate}[label=(\roman*)]
      \item Define atomic number and its relation to protons. (3)
      \item Define mass number and its components. (3)
      \item Explain the difference between atomic and mass numbers. (2)
    \end{enumerate}
  \end{questionblock}

  \begin{questionblock}[15]
    \item \textbf{(14 Marks)}

    (a) Describe the structure and trends in the periodic table. (6)
    \begin{enumerate}[label=(\roman*)]
      \item Define a period. (2)
      \item Define a group. (2)
      \item Describe a key trend across a period. (2)
    \end{enumerate}

    (b) Explain the predictive uses of the periodic table. (8)
    \begin{enumerate}[label=(\roman*)]
      \item How is element type identified? (3)
      \item How is reactivity predicted? (2)
      \item What can be deduced from an element's position? (3)
    \end{enumerate}
  \end{questionblock}

  \begin{questionblock}[15]
    \item \textbf{(14 Marks)}

    (a) Describe three main types of chemical bonds. (8)
    \begin{enumerate}[label=(\roman*)]
      \item Describe ionic bonding. (2)
      \item Explain covalent bonding. (3)
      \item What characterizes metallic bonding? (3)
    \end{enumerate}

    (b) Relate chemical bonding to substance properties. (6)
    \begin{enumerate}[label=(\roman*)]
      \item How does ionic bonding affect melting points? (2)
      \item What are typical properties of covalent compounds? (2)
      \item How does metallic bonding influence electrical conductivity? (2)
    \end{enumerate}
  \end{questionblock}

  \begin{questionblock}[15]
    \item \textbf{(14 Marks)}

    (a) Describe the atomic structure, including nucleus contents and electron arrangement. (8)
    \begin{enumerate}[label=(\roman*)]
      \item What is in the nucleus? (3)
      \item How are electrons arranged? (3)
      \item What is the atom's overall charge? (2)
    \end{enumerate}

    (b) Explain atomic number and mass number. (6)
    \begin{enumerate}[label=(\roman*)]
      \item Define atomic number. (2)
      \item Define mass number. (2)
      \item How do atomic and mass numbers differ? (2)
    \end{enumerate}
  \end{questionblock}

  \begin{questionblock}[15]
    \item \textbf{(14 Marks)}

    (a) Describe the structure of the periodic table. (8)
    \begin{enumerate}[label=(\roman*)]
      \item Define a period. (2)
      \item Define a group. (3)
      \item State a key trend across a period. (3)
    \end{enumerate}

    (b) Explain the predictive uses of the periodic table. (6)
    \begin{enumerate}[label=(\roman*)]
      \item How is an element's type identified? (2)
      \item How is reactivity predicted? (2)
      \item What can be deduced from an element's position? (2)
    \end{enumerate}
  \end{questionblock}

\end{enumerate}

% section C
\Needspace{16\baselineskip}
\centerheading{SECTION C (15 Marks)}
\vspace{-2.0em}
\begin{center}
  Answer only one (1) question in this section.
\end{center}
\vspace{-1.5em}

\begin{enumerate}[start=13]
  \begin{questionblock}[10]
    \item The atomic structure is a fundamental concept in chemistry. It is essential to understand its components and their significance. Explain the atomic structure, focusing on its components and the significance of atomic number and mass number. \textbf{(15 Marks)}
    \begin{enumerate}
      \item Describe the components of an atom and their arrangement. (5)
      \item Explain how the atomic number and mass number are related to the atomic structure. (10)
    \end{enumerate}
  \end{questionblock}

  \begin{questionblock}[10]
    \item The periodic table is a fundamental tool in chemistry that organizes elements based on their properties. Describe the structure and trends in the periodic table and explain its predictive uses. \textbf{(15 Marks)}
    \begin{enumerate}
      \item Describe the main features and trends observed in the periodic table. (7)
      \item How is the periodic table used to predict the properties and behavior of elements? (8)
    \end{enumerate}
  \end{questionblock}

\end{enumerate}

\end{document}
